\silentchapter{5 Det finst agentar}{det finst agentar}

\prop{5}{}{Verda rommar agentar som responderer på landskapet.}

Dei fire føregåande kapitla har skildra ein verd av former, krefter og landskap. Men dei har ikkje forklart kven eller kva som navigerer dette landskapet. Kreftene åleine produserer ikkje form; dei produserer eit felt av moglegheiter og restriksjonar. For at ei form skal realiserast, må noko eller nokon operere i dette feltet: oppdage gradienten, velje ein retning, utføre ein operasjon. Det er oppgåva til agenten.

Omgrepet «agent» vekkjer assosiasjonar til menneskelege designarar; til arkitekten ved teiknebrettet. Men agenten i denne traktaten er eit langt breiare omgrep. Materialet er ein agent. Marknaden er ein agent. Ein cellers bioelektriske nettverk er ein agent. Ein slimsopp som navigerer eit matlabyrint, er ein agent. Eit nevralt nettverk som genererer bilete, er ein agent. Dei deler alle same formelle eigenskap: dei responderer på ein gradient og justerer seg mot eit mål.

\section*{Kva er ein agent?}

\prop{5.1}{a}{Kvar klasse har minst eitt system som navigerer landskapet via negativ tilbakekopling.}

\prop{5.11}{d}{Ein \textbf{agent} er ein operator definert av trippelet: måltilstand, avstandsmåling, justeringsmekanisme. Agens krev berre organisering; verken medvit, intensjon eller nervesystem. Informasjonsflyt skil agentar frå passiv materie.}

Definisjonen er strengt. Ein agent er ikkje eit subjekt; han er eit system med tre eigenskapar: han har eit mål (ein tilstand han nærmar seg), han kan måle avstanden til målet, og han kan justere seg for å redusere den avstanden. Desse tre eigenskapane; mål, måling, justering; er alt som trengst. Definisjonen krev ikkje medvit. Ho krev ikkje intensjon. Ho krev ikkje eit nervesystem. Ho krev berre organisering; berre at systemet er sett opp slik at informasjonen om avstanden til målet faktisk påverkar kva systemet gjer neste gong.

\propcont{Formelt: Agent A := (g, d, \msym{δ}). Det finst ein Lyapunov-funksjon V slik at:}

\begin{center}\small
\begin{tabular}{l@{\hspace{4mm}}l}
V(g) = 0 & målet er nullpunktet \\
V(x) > 0 for x \msym{≠} g & alt anna er avstand \\
V er ikkje-aukande langs x\textsubscript{n+1} = \msym{δ}(x\textsubscript{n}) & systemet nærmar seg målet \\
\end{tabular}
\end{center}

Lyapunov-funksjonen gjev definisjonen matematisk presisjon. Eit system som oppfyller desse vilkåra, navigerer mot målet uavhengig av korleis det er implementert fysisk. Implementeringa kan vere biologisk (celler som justerer membranpotensial), mekanisk (ein termostat som justerer temperatur), algoritmisk (eit nevralt nettverk som justerer vekter) eller sosial (ein marknad som justerer prisar). Substratet er irrelevant; det er organiseringa som tel.

Ein rullande stein er ingen agent. Steinen rører seg, men rørsla manglar avstandsmåling. Steinen responderer ikkje på informasjon om kvar han er i høve til kvar han skal; han responderer utelukkande på tyngdekrafta i dette augeblikket. Det finst ingen Lyapunov-funksjon; ingen konvergens mot eit mål. Informasjonsflyt skil agentar frå passiv materie.

\prop{5.12}{o}{Agens er stratifisert. Termostaten og arkitekten skil seg berre i navigasjonskompleksitet.}

Dersom agentdefinisjonen ikkje krev medvit, følgjer det at agens kjem i gradar. Ein termostat er ein agent med éin dimensjon: temperatur. Ein arkitekt er ein agent med mange dimensjonar: lys, rom, struktur, materialkostnad, reguleringsplan, sosial kontekst. Dei er begge agentar; dei skil seg i kor mange aksar lyskjegla deira rommar.

Dette er ikkje ein metafor. Det er ein formell påstand om at same matematiske struktur; mål, måling, justering; opererer på tvers av alle desse systema. Skilnaden mellom termostaten og arkitekten er ikkje at den eine er ein «ekte» agent og den andre ikkje. Skilnaden er navigasjonskompleksitet: kor mange dimensjonar av formrommet agenten kan representere og justere for samstundes.

\section*{Navigasjon utan kart}

\prop{5.13}{t}{Agenten responderer på gradienten i landskapet, ikkje på forma. Han oppfinn ikkje målet; han oppdagar det. Navigasjon krev berre den lokale gradienten. Lokal kompetanse produserer global manifestasjon.}

Dette er ein av traktatens mest motintuitive påstandar. Me er vane med å tenkje at designaren har ein visjon; ein global plan; og at formgjevinga er prosessen med å realisere den planen. Men modellen seier noko anna. Agenten treng ikkje ein global plan. Han treng berre den lokale gradienten: i kva retning fell tilpassingslandskapet frå der han står no? Han tek eit steg i den retninga. Deretter måler han gradienten på nytt, og tek eit nytt steg.

Denne prosessen; gjenteken gradientnedstigning; produserer global form utan at nokon agent nokon gong hadde ein global representasjon av resultatet. Kvar agent opererer lokalt. Summen av dei lokale operasjonane er den globale forma. Lokal kompetanse produserer global manifestasjon.

Innsikta er ikkje at designarar ikkje tenkjer. Ho er at tenkinga ikkje treng å vere global for å produsere globale resultat. Ein stolmakar som justerer vinkelen på ein rygg inntil han «kjennest rett», følgjer ein lokal gradient i det fenomenologiske landskapet utan å ha eit kart over heile formrommet. Resultatet er likevel ein posisjon i formrommet som reflekterer alle dei seleksjonstrykka som verka på han i det augeblikket.

\section*{Den kognitive lyskjegla}

\prop{5.2}{d}{\textbf{Kognitiv lyskjegle} er delmengda av formrommet agenten kan representere og manipulere. Alt utanfor lyskjegla manglar kausal eksistens for agenten. Lyskjegler spenner frå cellulær millisekundrespons til tiårslang konsensus.}

Michael Levin og Chris Fields har formalisert omgrepet «kognitiv lyskjegle» som det romleg-temporale vindauget ein agent kan representere og påverke. Omgrepet er ein generalisering av det fysiske omgrepet «lyskjegle» frå relativitetsteorien: akkurat som det finst ein grense for kor langt lys kan reise i ein gjeven tidsperiode, finst det ei grense for kor mykje av formrommet ein agent kan representere og manipulere innanfor sin eigen tidsskala.

For ein celle er lyskjegla ein millisekundrespons på lokale biokjemiske signal. For ein handverkar er ho eit arbeidsrom definert av hendene, verktøyet og materialet framfor han. For ein marknad er ho ein tiårslang distribuert konsensusprosess som involverer millionar av aktørar. Alle tre er lyskjegler; dei skil seg i volum og tidsskala.

Alt utanfor lyskjegla manglar kausal eksistens for agenten. Ein stolmakar som ikkje veit at det finst eit marked for flatpakka møblar, kan ikkje navigere mot det. Ein celle som ikkje kan detektere vekstfaktor frå nabovevet, kan ikkje respondere på det. Grensa til lyskjegla er grensa til agentens verd.

\prop{5.21}{t}{Lyskjegla dikterer oppløysinga til formrommet. Agenten persiperer affordansar, ikkje former. Ein affordans er den lokale asymmetrien innanfor lyskjegla. Utviding aktualiserer latente affordansar; innsnevring deaktiverer reelle.}

Agenten ser ikkje former; han ser affordansar. Gibson sin innsikt var at persepsjon ikkje handlar om å oppfatte eigenskapar ved objekt, men om å oppdage kva operasjonar dei tilbyr. Ein stol vert ikkje persipiert som «eit objekt med fire bein og ein rygg»; han vert persipiert som «noko som tilbyr sitting». Affordansen er den lokale asymmetrien innanfor agentens lyskjegle.

Utviding av lyskjegla aktualiserer affordansar som var latente. Ein designar som lærer seg å arbeide med stål, ser plutseleg affordansar i materialet som ho ikkje såg før: tynne røyr kan bøyast utan å knekke; sveisar kan vere usynlege; flatevne kan kombinerast med styrke. Desse affordansane var der heile tida; det var lyskjegla som mangla dei.

Innsnevring av lyskjegla deaktiverer affordansar som var reelle. Ein designar som berre ser kostnad, ser ikkje materialets estetiske kvalitetar. Affordansane forsvinn ikkje; dei vert usynlege. Og forma som oppstår under den innsnevra lyskjegla, ber signaturen av det designaren ikkje såg.

\section*{Formgrammatikken}

\prop{5.3}{d}{Agenten manipulerer formgeometrien direkte. Konfigurasjonar lukka under union og snitt utgjer ein \textbf{formalgebra}; \textbf{formgrammatikken} utfører diskrete overgangar i denne geometrien. Reglar fyrer på innleiring; inndatageometrien er einaste premiss.}

Korleis navigerer agenten i formrommet? Gjennom grammatikkoperasjonar. Formgrammatikken; utvikla av George Stiny og James Gips i 1972; er eit formelt system der reglar opererer på former gjennom innleiring: dei gjenkjenner eit mønster i den noverande forma og erstattar det med eit anna.

\propcont{Ein formgrammatikk er eit trippel SG = (S, R, \msym{ω}) der R = \{r\textsubscript{i} : (a\textsubscript{i}, b\textsubscript{i})\}:}

\begin{center}\small
\begin{tabular}{l@{\hspace{4mm}}l}
\textbf{Regelapplikasjon} & s \msym{→} (s \msym{−} \msym{τ}(a\textsubscript{i})) + \msym{τ}(b\textsubscript{i}) \\[2pt]
\textbf{Språk} & L(SG) = \{s \msym{∈} S : \msym{ω} \msym{→}* s\} \\[2pt]
\textbf{Fyring} & regelen fyrer for kvar innleiring av a\textsubscript{i} i s \\
\textbf{Premiss} & berre noverande geometri; ikkje konstruksjonshistorie \\
\end{tabular}
\end{center}

Grammatikken er ein maskin for å utforske formrommet steg for steg. Kvart steg; kvar regelapplikasjon; tek agenten frå éin posisjon til ein nabo-posisjon. Mengda av alle posisjonar som er tilgjengelege gjennom éin operasjon utgjer agentens umiddelbare fridomsrom; naboskapen til den noverande forma.

Det vesentlege ved grammatikken er at han opererer utelukkande på den noverande forma. Konstruksjonshistoria er irrelevant. Regelen bryr seg ikkje om korleis forma kom til; berre om ho inneheld mønsteret som utløyser regelen. Innleiringa er ein eigenskap ved notida, ikkje ved fortida.

\prop{5.31}{t}{Grammatikkoperasjonen er ein C\msym{→}C-partisjon: ho utvidar tanken. Realisering er ein C\msym{→}K-operasjon: ho avgjer han. Grammatikken definerer moglegheitsrommet; landskapet fastset trajektorien. Agenten forløyser konfigurasjonen som mest radikalt reduserer lokal spenning.}

Her møtest dei tre tradisjonane. Grammatikkoperasjonen i Stinys forstand er ein C\msym{→}C-partisjon i Hatchuel og Weils forstand: kvar gong ein regel fyrer, vert tanken spesifisert vidare. Og realiseringa; det at prototypen vert bygd, testkøyrt, evaluert; er ein C\msym{→}K-operasjon: tanken vert avgjort. Grammatikken definerer kva rørsler som er moglege; landskapet avgjer kva rørsler som er verdifulle. Agenten opererer i skjeringspunktet.

\section*{Substratet som agent}

\prop{5.4}{d}{Substratet fungerer som ein null-ordens agent. Det vetoar konfigurasjonar og gjev resonans til andre gjennom ufravikelege avgrensingar.}

\prop{5.41}{i}{Biologisk vev navigerer via bioelektrisk polarisering; menneskeleg praksis via materiell friksjon; ein slimsopp via oscillerande samandragingar; ein språkmodell via vektoriell merksemd. Navigasjonsmekanikken er universell; kompetansen formelt ekvivalent.}

Substratet; det fysiske materialet som forma vert realisert i; er ikkje passiv mottakar. Det er sjølv ein agent, om enn av den enklaste typen. Bøketreet vetoar sprø bøyingar og gjev resonans til mjuke kurver. Stålet vetoar tjukkleik under ein viss terskel men opnar for langstrekte profil. Plasten vetoar ingenting av det treet vetoar, men innfører sine eigne forbod: UV-degradering, krimp, emissjonar.

Michael Levin har vist at dette gjeld heilt ned til det cellulære nivået. Celler navigerer i eit morforom; det bioelektriske mønsteret på cellemembranen kodar eit «målandlet» som vevet arbeider mot. Salamanderen som regenererer ein arm, navigerer mot eit mål; det opphavlege mønsteret; via ein Lyapunov-funksjon i det bioelektriske rommet. Slimsoppen \textit{Physarum polycephalum} finn det optimale rutenettet i eit matlabyrint via oscillerande samandragingar som konvergerer mot energiminima. Ein stor språkmodell navigerer det semantiske landskapet via vektoriell merksemd og produserer sekvensielle konfigurasjoner som konvergerer mot statistiske attraktorar.

Navigasjonsmekanikken er universell. Kompetansen til desse systema er formelt ekvivalent: han målast i evna til å minimere spenning langs dei lokale aksane i deira respektive lyskjegler. Salamanderen og språkmodellen utfører same formelle operasjon; berre substratet er ulikt.

\section*{Omsorg}

\prop{5.5}{t}{Navigasjonskompetanse akkumulerer seg som strukturert informasjon. Læring er den irreversible ekspansjonen av tilgjengeleg formrom.}

Kvar gong ein agent navigerer til ein ny posisjon i formrommet, utvidar han lyskjegla si. Posisjonen er erobra; ho er no ein del av K. Frå den nye posisjonen ser agenten nye affordansar som ikkje var synlege frå den førre. Læring er denne irreversible ekspansjonen av det tilgjengelege formrommet.

Ein handverkar som har laga tusen stolar, navigerer eit anna landskap enn ein som har laga éin. Ikkje fordi landskapet har endra seg, men fordi lyskjegla hans er utvida. Han ser gradientar som novisen ikkje ser. Han oppdagar kompromiss som novisen ikkje veit finst. Denne akkumulerte kompetansen er ikkje berre kunnskap; det er ein strukturell endring i agentens kapasitet til å representere formrommet.

\prop{5.6}{d}{\textbf{Omsorg} er mengda av aksar i formrommet agenten aktivt representerer og justerer for. Ein agent som held tre dimensjonar i lyskjegla, navigerer eit rikare landskap enn ein agent som held éin. Den fyrste oppdagar kompromiss den andre ikkje veit finst. Den fyrste produserer former som toler fleire trykk samstundes.}

No kjem me til det omgrepet som bind heile traktaten saman. Omsorg er ikkje eit kjenslemessig omgrep; det er eit formelt omgrep. Det refererer til mengda av aksar i formrommet; mengda av dimensjonar; som agenten aktivt held i lyskjegla si og justerer for. Ein agent med omsorg for éin dimensjon navigerer eit eindimensjonalt landskap. Ein agent med omsorg for ti dimensjonar navigerer eit tidimensjonalt landskap. Dei opererer i same formverda, men dei ser ulike landskap fordi lyskjeglene deira er ulike.

Kvifor kallar me det «omsorg»? Fordi det ikkje finst eit betre ord for kva det er å halde noko i merksemd og justere for det. Å ha omsorg for ein dimensjon er å sjå han; å la han verke på vala ein tek; å inkludere han i kompromisset. Å mangle omsorg for ein dimensjon er å vere blind for han; å overlevere han til krefter ingen representerer.

\prop{5.61}{t}{Omsorg er intelligensens innhald; intelligens er omsorgas form.}

\propcont{Ein agent utan omsorg for ein dimensjon er blind for landskapets gradient langs den aksen. Blindskapen er ikkje nøytral: dei usynlege trykka verkar framleis. Forma som oppstår under redusert omsorg ber signaturen av det agenten ikkje såg; ho sviktar der ho aldri vart prøvd. Omsorga kan ikkje supplerast i ettertid; den måtte vore der då forma vart navigert.}

Dette er traktatens kronjuvel. Omsorg og intelligens er ikkje to separate ting; dei er same ting sett frå to sider. Omsorg er kva intelligensen handlar *om*; intelligens er korleis omsorga er *organisert*. Ein agent som bryr seg om mange dimensjonar samstundes og som kan organisere omsorga slik at ho navigerer landskapet effektivt; det er ein intelligent agent.

Konsekvensen er radikal. Blindskapen er ikkje nøytral. Trykka langs den usynlege aksen verkar framleis; dei forsvinn ikkje fordi agenten ikkje ser dei. Forma som oppstår under ein agent med redusert lyskjegle; ein agent som berre bryr seg om kostnad, eller berre om estetikk, eller berre om effektivitet; ber signaturen av det agenten ikkje såg. Ho sviktar der ho aldri vart prøvd. Og svikten kan ikkje reparerast i ettertid; omsorga måtte ha vore der då forma vart navigert.

\prop{5.62}{t}{Å utvide lyskjegla er ein etisk operasjon: det er å ta fleire krefter med i kompromisset. Å snevra ho inn er å overlevere delar av formverda til krefter ingen representerer.}

Her vert formteori til etikk. Ikkje fordi traktaten feller verdidomar; ho gjer det ikkje. Men fordi strukturen i modellen viser at val av lyskjeglebreidde har konsekvensar som ikkje kan gjerast ugjorte. Å velje å ignorere materialets kompetanse er eit val med formelle konsekvensar. Å velje å ignorere brukarens affordansestruktur er eit val med formelle konsekvensar. Å snevra lyskjegla er å overlevere delar av formverda til krefter ingen representerer. Formene som oppstår i dei urepresenterte dimensjonane, er ikkje «designa»; dei er etterlatte.

\begin{overgang}
Agenten navigerer ikkje åleine. Kapittel 6 viser korleis fleire agentar, på ulike skalaer, med ulike lyskjegler, saman produserer den forma me til slutt observerer.
\end{overgang}
